%! Logarithmic Function Context and Data Modeling — Unit 2, Lesson 2.14 — Homework
\documentclass[10pt]{article}
\usepackage{precalculus-article}
\usepackage{precalculus-boxes}
\newcommand{\TallMath}[1]{$\displaystyle #1\rule[-1.4em]{0pt}{3.2em}$}

\begin{document}

\pageheader{Unit 2, Lesson 2.14}{Homework}
\namedateperiod

\vspace{0.1in}

\begin{enumerate}[leftmargin=1.8em, itemsep=14pt]

\item For each table, decide whether the data are better modeled by a \textbf{logarithmic}
      or an \textbf{exponential} function. Circle your choice and give a one-sentence
      justification by examining how inputs change for equal output increments (for log)
      or how outputs change for equal input increments (for exponential).

\vspace{6pt}
\begin{center}
\renewcommand{\arraystretch}{1.8}
\begin{tabular}{cc}
  \begin{tabular}{|c|c|}
  \hline
  \rowcolor{sky}
  $x$ & $y$ \\
  \hline
  $1$ & $0$ \\
  \hline
  $3$ & $1$ \\
  \hline
  $9$ & $2$ \\
  \hline
  $27$ & $3$ \\
  \hline
  \end{tabular}
  \qquad\qquad &
  \begin{tabular}{|c|c|}
  \hline
  \rowcolor{sky}
  $x$ & $y$ \\
  \hline
  $0$ & $2$ \\
  \hline
  $1$ & $6$ \\
  \hline
  $2$ & $18$ \\
  \hline
  $3$ & $54$ \\
  \hline
  \end{tabular}
\end{tabular}
\end{center}

\vspace{4pt}
\textbf{Table A:} \quad Logarithmic \quad Exponential \quad because \writeline

\vspace{4pt}
\textbf{Table B:} \quad Logarithmic \quad Exponential \quad because \writeline

\item A data set contains the points $f(1) = 0$ and $f(e^3) = 6$. Assume
      $f(x) = a\ln(x)$.

\begin{enumerate}[leftmargin=1.5em, itemsep=8pt, label=\alph*.]
  \item From $f(1) = 0$: verify that $\ln(1) = 0$, so any value of $a$ satisfies this.
        Use $f(e^3) = 6$ to find $a$.

  \vspace{4pt}
  $a\ln(e^3) = 6 \;\Rightarrow\; a \cdot \blank{1.5cm} = 6 \;\Rightarrow\; a = $ \blank{2cm}

  \item Write the complete model.

  \vspace{4pt}
  $f(x) = $ \blank{4cm}

  \item Predict $f(e^5)$.

  \vspace{4pt}
  $f(e^5) = $ \blank{5cm}
\end{enumerate}

\item The table below shows habitat area $A$ (hectares) and species count $S$.

\vspace{6pt}
\begin{center}
\renewcommand{\arraystretch}{1.8}
\begin{tabular}{|c|c|c|c|c|}
\hline
\rowcolor{sky}
$A$ & $1$ & $10$ & $100$ & $1000$ \\
\hline
$S$ & $3$ & $6$  & $9$   & $12$ \\
\hline
\end{tabular}
\end{center}

\begin{enumerate}[leftmargin=1.5em, itemsep=8pt, label=\alph*.]
  \item Use the two points $(1,\,3)$ and $(10,\,6)$ to write a system of equations
        for $S(A) = a\log_{10}(A) + k$.

  \vspace{4pt}
  Equation 1: \blank{6cm} \quad Equation 2: \blank{6cm}

  \item Solve for $k$ and $a$.

  \vspace{4pt}
  $k = $ \blank{2cm} \qquad $a = $ \blank{2cm}

  \item Write the complete model and predict $S(10000)$.

  \vspace{4pt}
  $S(A) = $ \blank{4cm} \qquad $S(10000) = $ \blank{3cm}
\end{enumerate}

\item Let $f(x) = 2\log_2(x - 1) + 3$.

\begin{enumerate}[leftmargin=1.5em, itemsep=8pt, label=\alph*.]
  \item State the domain of $f$.

  \vspace{4pt}
  Domain: \blank{5cm}

  \item Evaluate $f(5)$.

  \vspace{4pt}
  $f(5) = 2\log_2(5-1)+3 = 2\log_2(\blank{1cm})+3 = 2\cdot\blank{1cm}+3 = $ \blank{2cm}

  \item Solve $f(x) = 7$ for $x$. Show all steps.

  \vspace{4pt}
  \writelines{3}
  $x = $ \blank{3cm}
\end{enumerate}

\item \textbf{AP-style multiple choice.} Which situation is best modeled by a
      \emph{logarithmic} function?

\vspace{6pt}
\begin{enumerate}[leftmargin=2em, itemsep=4pt, label=(\Alph*)]
  \item A bank account balance that grows by 5\% each year.
  \item A child's height from birth to adulthood: grows rapidly at first, then slows
        significantly.
  \item The distance a car travels at constant speed, which is proportional to time.
  \item A bacterial colony that triples in size every 3 hours.
\end{enumerate}

\item \textbf{Extension.} Using the change-of-base property $\log_b x = \dfrac{\ln x}{\ln b}$,
      explain why $f(x) = a\ln(x)$ and $g(x) = a\log_{10}(x)$ always produce models that
      are vertical dilations of each other. Write the exact dilation factor that converts
      one into the other.

\writelines{3}

\end{enumerate}

\vspace{0.12in}

\begin{reflectionbox}
What is one thing from today's lesson that you feel confident about, and one question you
still have about constructing or using logarithmic models?

\writelines{2}
\end{reflectionbox}

\end{document}
